\section{Jenis dan Pendekatan Penelitian}

Penelitian ini menggunakan pendekatan \textit{Design Science Research} (DSR) yang diperkuat dengan evaluasi eksperimental berbasis simulasi. Pendekatan DSR digunakan karena penelitian ini tidak hanya bertujuan menganalisis permasalahan kontrol akses pada lingkungan IoT, tetapi juga merancang artefak berupa mekanisme \textit{GWO-optimized ABAC} untuk mengoptimalkan \textit{cache} atribut dan \textit{attribute freshness} pada ABAC berbasis \textit{edge}. Artefak yang dirancang mencakup arsitektur ABAC berbasis \textit{edge}, mekanisme \textit{cache} atribut lokal pada \textit{gateway}, formulasi optimasi \textit{cacheable attribute}, nilai TTL per atribut, ukuran \textit{cache}, serta fungsi \textit{fitness} yang mempertimbangkan aspek performa dan kualitas keputusan akses.

Secara metodologis, penelitian ini bersifat kuantitatif eksperimental karena evaluasi dilakukan melalui pengukuran metrik numerik pada beberapa konfigurasi pengujian. Konfigurasi yang dibandingkan meliputi \textit{cloud-only ABAC, edge-fresh ABAC, static cache ABAC, adaptive TTL ABAC, random search ABAC}, dan \textit{GWO-optimized ABAC}. Perbandingan tersebut digunakan untuk memisahkan pengaruh penempatan otorisasi di \textit{edge}, penggunaan \textit{cache} atribut, strategi \textit{freshness}, dan optimasi metaheuristik terhadap performa serta kualitas keputusan otorisasi.

Eksperimen dilakukan melalui simulasi berbasis iFogSim2 dan implementasi Java untuk merepresentasikan lingkungan \textit{edge} IoT dengan topologi \textit{cloud–proxy–gateway–endpoint}. Simulasi dipilih karena memungkinkan pengaturan skenario secara terkontrol, termasuk jumlah \textit{gateway} dan \textit{endpoint}, dinamika \textit{workload}, perubahan atribut, ketersediaan PIP, serta kejadian revokasi. Dengan demikian, setiap konfigurasi dapat diuji pada kondisi yang sama sehingga hasil perbandingan menjadi lebih objektif dan dapat direproduksi.

Evaluasi penelitian difokuskan pada dua aspek utama, yaitu efisiensi performa dan kualitas keputusan akses. Aspek performa diukur melalui \textit{average latency, p95 latency,} jumlah \textit{PIP calls, cache hits, CPU time,} dan penggunaan memori. Sementara itu, kualitas keputusan akses diukur melalui \textit{false allow rate, false deny rate, stale cache hit rate,} dan \textit{revocation allow rate}. Dengan pendekatan tersebut, penelitian ini tidak hanya menilai konfigurasi yang paling cepat, tetapi juga menilai konfigurasi yang paling mampu menjaga keseimbangan antara efisiensi sistem, kesegaran atribut, dan keamanan keputusan otorisasi pada lingkungan IoT \textit{edge} yang dinamis.